\section{The retained family and its intrinsic classification}\label{fam:section:A}

Throughout this section $\mathcal P=\{c:\operatorname{Im}c<-M\}$, and $X_c$ means the actual glued space with $v_1=\varepsilon=(1,2,-4,0)^t$, $v_2=-\varepsilon'=(-1,-3,3,0)^t$ and the untwisted cusp collar. We write $d=\widehat\delta$, $e=(0,1)^t$, $E=d\gamma$, and use column vectors. Our peripheral generators are clockwise: $s_j(g_jz)=e^{-2\pi i/m_j}s_j(z)$ and $g_1g_2g_0=1$. Thus $s(g_0z)=s(z)-1$ at the cusp. All period comparisons use
\[
 A_1=\begin{pmatrix}1&0&0&0\\6&0&1&0\\-6&-1&-1&0\\-2&1&0&1\end{pmatrix},\quad
 A_2=\begin{pmatrix}1&0&0&0\\0&0&-1&0\\-6&1&0&0\\3&0&1&1\end{pmatrix},\qquad
 R_g(z)\Pi_c(z)=\Pi_c(gz)A_g.
\]
Here $R_1=\left(\begin{smallmatrix}-1/\tau&0\\(1-\mu)/\tau&1\end{smallmatrix}\right)$, $R_2=\left(\begin{smallmatrix}1/\tau&0\\-\mu/\tau&1\end{smallmatrix}\right)$, and $R_0=I$. The distinction between source and target markings in each equation below is retained by its displayed period equation.

\subsection{Canonical sections and the finite covers}\label{fam:section:A.1}

\begin{proposition}\label{fam:src:A.1.1}\label{fam:st:A.1.1}
For the constructed holomorphic map $f_c:X_c\to\mathbb P^1$, and also for every admissible choice of its affine translations and cusp collar, we have
\[
 \omega_{X_c}\simeq\mathcal O_{X_c}(-2S_2),\qquad
 f_{c*}\mathcal O_{X_c}(kS_2)=\mathcal O_B(\lfloor k/4\rfloor p_2)\quad(k\in\mathbb Z).
\]
The latter equality is an equality of subsheaves of meromorphic functions under pullback. The anticanonical section ring is $\mathbb C[U,V]$, with $\deg U=1$, $\deg V=2$. The complete degree-two map is $f_c$ followed by a projective coordinate isomorphism of its base. Every biholomorphism between two such threefolds is over the identity of $B$.
\end{proposition}
\Uses{the hypotheses, notation and definitions stated here; \citet[Theorem~9.1]{AC2026}; \cref{app:rederive:fam-src-A.1.1}}


The canonical pencil is due to \citet[Theorem~9.1]{AC2026}. A complete re-derivation is given in Appendix~\ref{app:rederive:fam-src-A.1.1}.


\begin{proposition}\label{fam:src:A.1.2}\label{fam:st:A.1.2}
The normalized base change $N_j\times_{D_j}\Delta_j\to\Delta_j$, $t_j=s_j^{m_j}$, is the original smooth torus family $J'_j$. The maps $E_1\times F_1\to S_1$ and $E_2\times F_2\to S_2$ described below are finite etale covers of degrees $9$ and $8$.
\end{proposition}
\Uses{the hypotheses, notation and definitions stated here; the explicit constructions and calculations in the proof}

\begin{proof}
On $J'_j$, the deck generator is
$(s,\zeta)\mapsto(\zeta_js,R_j(s)\zeta+\Pi_c(g_js)v_j/m_j)$. It has exact order $m_j$. Its $r$th power changes the invariant circle coordinate $\gamma$ by $r\gamma(v_j)/m_j$; this is nonintegral for $0<r<m_j$, so the action is free. The quotient map is therefore a local biholomorphism. Its local base function is $s^{m_j}$ and its fibre multiplicity is exactly $m_j$. The map
\[
 J'_j\longrightarrow N_j\times_{D_j}\Delta_j,\qquad y\longmapsto(q_j(y),s_j(y))
\]
is finite, is an isomorphism over $s\ne0$, and has normal source. Locally at the central fibre the base-change ring has equation $z^{m_j}-s^{m_j}=0$. Its normalization is the disjoint union of branches $z=\zeta_j^rs$. The displayed map gives those branches in the covering charts. It is thus the normalization everywhere, with its specified root coordinate.

Put $\rho=e^{\pi i/3}$, $\mu_1=(2-\rho)/3$, $\mu_2=(1-i)/2$, $\kappa_1=\beta_c(z_1)+2\mu_1$, and $\kappa_2=\beta_c(z_2)+3\mu_2$. The fixed-point transformation laws for $\mu$ give the two displayed values directly. We have $\operatorname{Im}\kappa_j=D_c(z_j)<0$. Define
\[
 E_j=\mathbb C/(\mathbb Z+\kappa_j\mathbb Z),\quad
 F_1=\mathbb C/(\mathbb Z(\rho-2)-3\mathbb Z),\quad
 F_2=\mathbb C/(\mathbb Z(i-1)-2\mathbb Z).
\]
Let $e_1=(0,1,-2,0)^t,f_1=(0,0,-3,1)^t,e_2=(0,1,-1,0)^t,f_2=(0,0,-2,1)^t$. Direct solution of the fixed equations and inspection of the columns give
\[
 \Lambda^{A_1}=\mathbb Z\varepsilon+\mathbb Zd,\quad
 \Lambda^{A_2}=\mathbb Z\varepsilon'+\mathbb Zd,\quad
 (A_j-I)\Lambda=\mathbb Ze_j+\mathbb Zf_j.
\]
The image lattices are saturated: they are the integer solutions of $\gamma=0,2u+w+3\delta=0$ and $\gamma=0,u+w+2\delta=0$, respectively. The determinants of the fixed-plus-moving basis matrices are $-3,-2$.

The complex-linear maps $L_1(\xi,\eta)=(\eta,\xi-\eta/3)$ and $L_2(\xi,\eta)=(\eta,\xi-\eta/2)$ carry the product periods to those fixed and moving sublattices. They induce isogenies to the original central tori, with kernel generated respectively by $(2/3,\rho)$ and $(1/2,i)$, of exact orders three and two. For example $\widehat u=2d/3+e_1-2f_1/3=d/2+e_2-f_2/2$ gives both kernel representatives.

The affine deck maps lift to
\[
 (\xi,\eta)\longmapsto(\xi+\kappa_1/3,(\rho-1)\eta),\qquad
 (\xi,\eta)\longmapsto(\xi-\kappa_2/4,-i\eta).
\]
Substitution in $L_j$ gives exactly $R_jL_j+\Pi_cv_j/m_j$. They commute with their kernel translations: the moving differences are periods, since $(\rho-2)\rho=-1-\rho=-(\rho-2)-3$ can equally be checked in the lattice basis, and $(-i-1)i=1-i=-(i-1)$. More directly these assertions follow by the two-dimensional period matrices of the indicated rotations. A nonidentity element of either resulting product group translates the $E_j$ coordinate by $2a/3+b\kappa_1/3$ or $a/2-b\kappa_2/4$, respectively. Real independence of $1,\kappa_j$ proves that this translation vanishes only for the identity element. Hence the groups $C_3\times C_3$ and $C_2\times C_4$ act freely and have quotient precisely $S_j$.
\end{proof}

\subsection{The integral obstruction}\label{fam:section:A.2}

\begin{proposition}\label{fam:src:A.2.1}\label{fam:st:A.2.1}
The rational centralizer of $(A_1,A_2)$ is $\mathbb QI+\mathbb QE$. Every biholomorphism $X_c\to X_{c'}$ has $c'-c=n\in2\mathbb Z$, complex-linear fibre part $I_2$, and source-to-target lattice map $I-nE$. In the full fixed lattices,
\[
 \Lambda^{A_1}/N_1\Lambda\simeq\mathbb Z/3,\qquad
 \Lambda^{A_2}/N_2\Lambda\simeq\mathbb Z/4\oplus\mathbb Z/2,
\]
with $d$ generating the second factor of the latter group and vanishing in the former.
\end{proposition}
\Uses{the hypotheses, notation and definitions stated here; \cref{fam:src:A.1.2}}

\begin{proof}
The common invariant vector space is $\mathbb Qd$, and the common invariant covector space is $\mathbb Q\gamma$; this follows by intersecting the fixed bases in \ref{fam:src:A.1.2} and the rows $\langle\gamma,2u+w+3\delta\rangle$, $\langle\gamma,u+w+2\delta\rangle$. A commuting endomorphism preserves $\mathbb Qd\subset\ker\gamma$. Its induced matrix on $\ker\gamma/\mathbb Qd$ commutes with
$Q_1=\left(\begin{smallmatrix}0&1\\-1&-1\end{smallmatrix}\right)$ and $Q_2=\left(\begin{smallmatrix}0&-1\\1&0\end{smallmatrix}\right)$. Commutation with $Q_2$ makes it $\left(\begin{smallmatrix}b&r\\-r&b\end{smallmatrix}\right)$, and the $Q_1$ equation sets $r=0$. Write the remaining columns as
$H\widehat\gamma=(a,x,y,k)^t$, $H\widehat u=b\widehat u+u_0d$, $H\widehat w=b\widehat w+w_0d$, $Hd=d_0d$. The equations on $\widehat u,\widehat w$ give $d_0-w_0=b+u_0$, $u_0=2w_0$ and $u_0=w_0$, hence $u_0=w_0=0,d_0=b$. The first-column equations give $x=2(a-b),y=-4(a-b),x=-y$. Thus $a=b,x=y=0$, proving the assertion. Conversely $A_jE=EA_j=E$.

A fibre biholomorphism lifts to an affine map on $\mathbb C^2$: the derivative of a lifted map is lattice periodic, descends to a holomorphic function on a compact torus, and is constant. Its lattice map is locally constant and commutes with both monodromies. Write it as $H=aI+bE$; integrality and $\det H=a^4=1$ give $a=\pm1,b\in\mathbb Z$. The equation $L\Pi_c=\Pi_{c'}H$, in columns three and four, gives $L=aI_2$; column one gives $a(c'-c)+b=0$. Thus $n=c'-c\in\mathbb Z$.

The normalized base-change construction in \ref{fam:src:A.1.2} is functorial for biholomorphisms over the base: the base-changed isomorphism transports integral elements of the total quotient ring to integral elements, and its inverse does the same. Its lift is over the identity of the root disc. It commutes with the prescribed deck generator because both compositions lift the same map $(x,s)\mapsto(F(x),\zeta_js)$ on the unnormalized base change. This proves exact commutation, rather than commutation after an unchosen power of the generator.

Lift the translation of the resulting affine torus map holomorphically to a vector $B_j(s)$ on the simply connected root disc. Its deck equation, expressed in the source periods, is
\[
 B_j(g_jz)-R_j(z)B_j(z)=\Pi_c(g_jz)
 \left(\frac{(1-a)v_j+n\gamma(v_j)d}{m_j}+\ell_j\right),\qquad\ell_j\in\Lambda.
\]
The integral vector is constant by continuity and real invertibility of the period matrix. Iterating the equation $m_j$ times gives
\[
 N_j\ell_j=-(1-a)v_j-n\gamma(v_j)d.
\]
Exact multiplication yields
\[
 N_1=\begin{pmatrix}3&0&0&0\\6&0&0&0\\-12&0&0&0\\0&2&1&3\end{pmatrix},\quad
 N_2=\begin{pmatrix}4&0&0&0\\12&0&0&0\\-12&0&0&0\\0&2&2&4\end{pmatrix}.
\]
In fixed bases their images are $3\mathbb Z\varepsilon+\mathbb Zd$ and $4\mathbb Z\varepsilon'+2\mathbb Zd$. These explicit generators prove the asserted quotient groups. Applying $\gamma$ to the norm equation excludes $a=-1$ already at $m_1=3$, since $3\gamma(\ell_1)=-2$. For $a=1$, the second norm equation says exactly $2\mid n$; the first imposes no condition. Program \texttt{code/ab\_exact.py} verifies the matrices, Smith forms, centralizer and membership independently by Hermite-form comparisons. All its computations are over $\mathbb Z$ or $\mathbb Q$.
\end{proof}

\subsection{Extension and coherent descent}\label{fam:section:A.3}

\begin{proposition}\label{fam:src:A.3.1}\label{fam:st:A.3.1}
For each $n\in2\mathbb Z$ there is a biholomorphism $F_{n,c}:X_c\to X_{c+n}$, holomorphic jointly in $c$, with $F_{2h,c+2k}F_{2k,c}=F_{2(h+k),c}$. The quotient family over $0<|w|<e^{-\pi M}$, $w=e^{-\pi ic}$, is proper and holomorphic, with pairwise non-biholomorphic fibres.
\end{proposition}
\Uses{the hypotheses, notation and definitions stated here; \cref{fam:src:A.2.1}}

\begin{proof}
Integer shifts satisfy $\Pi_{c+n}=\Pi_c(I+nE)$, so the middle lattices agree. Set
\[
 r_1=-\tfrac n3(\widehat u+\widehat w),\quad
 r_2=\tfrac n4(\widehat u+\widehat w),\quad
 \ell_1=-n\widehat w,\quad \ell_2=\tfrac n2\widehat u,\quad b_j=\Pi_cr_j.
\]
The four-entry equations are $(I-A_1)r_1=(0,0,-n,n/3)^t$ and $(I-A_2)r_2=(0,n/2,0,-n/4)^t$. Consequently $B_j=b_j$ solves the preceding deck equation for $a=1$. The map $(z,\zeta)\mapsto(z,\zeta+b_j(z))$ extends on the entire root family and descends, together with its inverse, to each finite filling.

Its middle translation on the punctured overlap is
\[
 a_j=b_j+n\gamma(v_j)\frac{\log s_j}{2\pi i}e.
\]
This follows from the actual overlap $h_{j,c}(\zeta)=\zeta+\sigma_{j,c}$, where $\sigma_{j,c}=(\log s_j)/(2\pi i)\,\Pi_cv_j$: precisely $\sigma_{j,c+n}-\sigma_{j,c}=n\gamma(v_j)(\log s_j)/(2\pi i)e$. Changes of logarithm add periods. Since the chosen generator changes $(\log s_j)/(2\pi i)$ by $-1/m_j$, we have $a_j(g_jz)=R_ja_j(z)+\Pi_c(g_jz)\ell_j$.

On the universal cover of $B^\circ$, extend $\ell_1,\ell_2$ by the integral cocycle law $\ell_{gh}=\ell_g+A_g\ell_h$. Freeness of $\pi_1(B^\circ)$ makes this well defined. At $g_0=(g_1g_2)^{-1}$ it gives
\[
 \ell_0=-A_0(\ell_1+A_1\ell_2)=\tfrac{3n}2\widehat w-\tfrac n2d.
\]
Put $q_n=(-3n/2,n/2)^t$ and $a_0=sq_n$. Then $a_0(g_0z)-a_0(z)=\Pi_c\ell_0$. In cusp torus coordinates its map is $(x,t_c)\mapsto(t_c^{q_n}x,t_c)$. The integral linear map $(y,h)\mapsto(y+hq_n,h)$ of the fan lattice translates each height-one triangle by an integer vector. It therefore permutes every cone of the prescribed periodic triangulation. The map and its inverse extend across all strata, including both triple points and all three double curves after quotient. Moreover the cusp deck multipliers are identical at $c$ and $c+n$, since $C_{c+n}\lambda-C_c\lambda=(0,n\lambda_1)^t$ has integral exponential. Thus the extended map descends with all its nonconstant units retained.

We supply the global analytic correction. Quotient the trivial affine vector bundle on the universal cover by
$(z,y)\mapsto(gz,R_g(z)y+\Pi_c(gz)\ell_g)$. Its action law is exactly the cocycle law by period covariance. Specify at the finite points sections of the form $a_j+\eta_j$, with $\eta_j$ holomorphic equivariant on the full root disc; at the cusp specify $a_0+\eta_0(t_c)$ with $\eta_0$ holomorphic on the full disc. These prescriptions define a locally nonempty torsor under a rank-two locally free sheaf $\mathcal V$ on $B$.

Projection to the first coordinate gives
\[
 0\longrightarrow\mathcal O_Be\longrightarrow\mathcal V
 \longrightarrow\mathcal O_B(-p_0)\longrightarrow0.
\]
Indeed its factors are $-1/\tau,1/\tau,1$. Multiplication by $\Theta$ identifies this first-coordinate sheaf with holomorphic functions at ordinary and finite points and functions vanishing at $p_0$: the finite vanishing orders are $2,1$, and the cusp pole is one. At a finite stalk, any equivariant first component has an equivariant vector lift by averaging a preliminary lift $(u,0)$ over the finite group. This proves stalkwise surjectivity and local freeness, including the off-diagonal entries of $R_j$.

Both $H^1(\mathbb P^1,\mathcal O)$ and $H^1(\mathbb P^1,\mathcal O(-1))$ vanish. To make the analytic gluing explicit, split a smooth cocycle with a partition of unity, solve $\bar\partial$ locally by the Cauchy integral, and compare the two resulting solutions on an annulus. A Laurent series $\sum a_nt^n$ is the difference of its nonnegative and negative parts for $\mathcal O$; for $\mathcal O(-1)$ use the second frame $t^{-1}$, so the negative terms become $t^{-1}(-\sum_{n<0}a_nt^{n+1})$. Both series converge on their respective discs and leave no exponent unsolved. This kills the cocycle for any refinement. The exact sequence gives $H^1(B,\mathcal V)=0$. Project the torsor cocycle to $\mathcal O(-1)$, solve it, lift the solutions, and solve the remaining scalar cocycle. We obtain one global section $a$, with $a=a_j+\eta_j$ at finite points and $a=sq_n+\eta_0$ at the cusp.

The three actual maps are now
\[
 \zeta\mapsto\zeta+a,\qquad
 \zeta\mapsto\zeta+b_j+\eta_j,\qquad
 (x,t_c)\mapsto(e^{2\pi i\eta_0(t_c)}t_c^{q_n}x,t_c).
\]
The finite overlap identity is the equality of sums
$\sigma_{j,c}+a=\sigma_{j,c+n}+b_j+\eta_j$. The cusp overlap identity is $e^{2\pi ia}=t_c^{q_n}e^{2\pi i\eta_0}$. The holomorphic-unit factors act on every toric chart. The local inverses satisfy the same identities; hence the maps glue to a biholomorphism.

For coherence, perform the additive solution only for $n=2$ and multiply that solution by $k$ for $n=2k$. The cocycles have $\gamma(\ell_g)=0$, so their affine terms $\Pi_c\ell_g$ are independent of $c$. The reference translations, $\mathcal V$, and the correction can all be chosen independent of $c$. Adding the translations and multiplying the toric factors proves the stated cocycle identity exactly.

Finally the original gluing constructs a holomorphic family over any parameter disc compactly contained in $\mathcal P$. In ordinary and finite charts, real invertibility of $\Pi_c$ is uniform on compact parameter sets, so a fixed compact real parallelepiped proves properness and the lattice quotient is holomorphic. At the cusp write $Z=sB_0+C(t_c,c)$ with $B_0=\left(\begin{smallmatrix}0&1\\-1&0\end{smallmatrix}\right)$. On a compact parameter disc $C$ is uniformly bounded; shrink the cusp radius until
$\|(-2\pi\operatorname{Im}C)/\log|t_c|\|<1/2$. The effective logarithmic lattice is then within $1/2$ of the unimodular $B_0$, uniformly. Reduce its logarithmic valuation modulo the effective logarithmic lattice to a bounded parallelogram. Choose a height-one triangle containing that reduced valuation. Its barycentric coordinates are nonnegative and sum to one; the corresponding three unimodular cone characters therefore have modulus at most one. This puts the point in the closed unit polydisc of that cone. Only finitely many triangles meet the fixed bounded parallelogram, uniformly on the parameter compactum. Boundary orbits have the same representatives by translating their height-one components to these triangles. Thus every orbit has a representative in this finite union of cone polydiscs. Those closed polydiscs times a compact parameter set form a compact surjective representative set. This proves properness at the cusp and proper discontinuity by the same uniform lower bound on nonzero lattice translations. The quotient charts are open subsets of a fixed smooth toric variety times the parameter disc, so projection to the parameter is a submersion. The explicit overlap maps are jointly holomorphic. These local families agree after shrinking cusp radii, and glue over $\mathcal P$. The coherent $2\mathbb Z$ action is properly discontinuous because its action on $\mathcal P$ is. Local inverse logarithms for $w=e^{-\pi ic}$ identify the quotient family with these proper local families. \ref{fam:src:A.2.1} gives pairwise non-biholomorphic fibres.
\end{proof}

\subsection{Albanese invariants and the derivative of the deformation}\label{fam:section:A.4}

\begin{proposition}\label{fam:src:A.4.1}\label{fam:st:A.4.1}
With $\kappa_j(c)=\kappa_j^{\rm part}+c$, the Albanese elliptic curves of the reduced finite fibres are
\[
 \operatorname{Alb}(S_1)=\mathbb C/(\mathbb Z+\kappa_1(c)\mathbb Z),\qquad
 \operatorname{Alb}(S_2)=\mathbb C/(\mathbb Z+\tfrac12\kappa_2(c)\mathbb Z).
\]
Writing $J(q)=q^{-1}+744+196884q+\cdots$, their invariants are
\[
 j_1(w)=J(e^{-2\pi i\kappa_1^{\rm part}}w^2),\qquad
 j_2(w)=J(e^{-\pi i\kappa_2^{\rm part}}w).
\]
These are convergent formulas because the displayed nome moduli are below one. Outside the explicitly defined countable exceptional set in the proof, equality of both Albanese invariants forces $c'-c\in2\mathbb Z$. Thus the generic Albanese invariants recover the parameter modulo its exact period subgroup, independently of the centralizer calculation.
\end{proposition}
\Uses{the hypotheses, notation and definitions stated here; \cref{fam:src:A.1.2}; \cref{fam:src:A.2.1}}

\begin{proof}
In the product covers of \ref{fam:src:A.1.2} the invariant holomorphic one-forms are exactly the multiples of $d\xi$: the moving factor is multiplied by a primitive third or fourth root. Integration thus identifies the Albanese with the quotient of the fixed line by the fixed-coordinate projections of every deck translation. These are, respectively,
\[
 \mathbb Z+\kappa_1\mathbb Z+\tfrac23\mathbb Z+\tfrac{\kappa_1}3\mathbb Z
   =\tfrac13\mathbb Z+\tfrac{\kappa_1}3\mathbb Z,
\quad
 \mathbb Z+\kappa_2\mathbb Z+\tfrac12\mathbb Z+\tfrac{\kappa_2}4\mathbb Z
   =\tfrac12\mathbb Z+\tfrac{\kappa_2}4\mathbb Z.
\]
Multiplication by three and two gives the stated maps of complex tori. In particular the Albanese computation includes the kernel of the product isogeny and the affine quotient, rather than just its fixed eigenspace. Using the positive-imaginary representatives $-\kappa_1$ and $-\kappa_2/2$ in the standard modular $j$-series gives the formulas. At sufficiently small nome $J$ is injective: $1/J(q)=q+O(q^2)$ has nonzero derivative at zero. Thus $j_2$ alone distinguishes parameters within a sufficiently small punctured subdisc when both parameters lie there. For the full generic assertion put $x=-\kappa_2(c)/2$ and $d_*=\kappa_1^{\rm part}-\kappa_2^{\rm part}$, so $-\kappa_1(c)=2x-d_*$. The independently certified interval in \texttt{cd\_special\_intervals.py} gives $-0.2<\operatorname{Im}d_*<-0.1$; in particular it is nonzero. Equality of the two elliptic invariants at $c,c'$ gives matrices $M,N\in\operatorname{PSL}_2(\mathbb Z)$ with $x'=M(x)$ and
\[
 N(2x-d_*)=2M(x)-d_*.
\]
Define the exceptional set to be the $c\in\mathcal P$ obtained from roots of these equations for all pairs for which the two rational functions are not identical. Each equation has at most two roots after clearing denominators, so this is a countable, explicitly specified set. For a pair giving an identity, let $C(x)=2x-d_*$. Then $N=CMC^{-1}$. The real matrix $M$ preserves both the real circle $\mathbb R\cup\{\infty\}$ and $C^{-1}(\mathbb R\cup\{\infty\})$, the parallel line with imaginary part $\operatorname{Im}d_*/2$ together with infinity. Their unique intersection is infinity, so $M$ fixes infinity. An element of $\operatorname{PSL}_2(\mathbb Z)$ fixing infinity is $x\mapsto x+b$, $b\in\mathbb Z$, and then $N(y)=y+2b$. It follows that $c'-c=-2b$. This proves the generic necessity with its actual exceptional set; \ref{fam:src:A.2.1} proves the necessity also on that set.
\end{proof}

\begin{proposition}\label{fam:src:A.4.2}\label{fam:st:A.4.2}
For every $c\in\mathcal P$, the Kodaira--Spencer class of $\partial/\partial c$ in $H^1(X_c,T_{X_c})$ is nonzero. On any fixed ordinary fibre, put $T=\operatorname{Im}\tau$, $m=\operatorname{Im}\mu$, $a=m/T$, $D=\operatorname{Im}\beta_c-6m^2/T$. In the convention $\mathrm{KS}(\partial_c)=[\bar\partial h]$ for a smooth $(1,0)$ lift $h$, its restricted torus-deformation representative is
\[
 -\frac{1}{2iD}\,\partial_{\zeta_2}\otimes
       (d\bar\zeta_2-a\,d\bar\zeta_1).
\]
\end{proposition}
\Uses{the hypotheses, notation and definitions stated here; \cref{fam:src:A.1.1}}

\begin{proof}
Retain the same real marking $r\in\mathbb R^4$ and set $\zeta=\Pi_cr$. The two imaginary-part equations give
$r_1=(\operatorname{Im}\zeta_2-a\operatorname{Im}\zeta_1)/D$. Thus at $c+\delta$ the marked real-linear identification is
\[
 w=\zeta+\frac{\delta}{2iD}P(\zeta-\bar\zeta),\qquad
 P=\begin{pmatrix}0&0\\-a&1\end{pmatrix},\qquad P^2=P.
\]
It is exactly $\Pi_{c+\delta}\Pi_c^{-1}$ as a real map, so it sends every marked lattice vector to the corresponding vector. Its determinant is $(D+\operatorname{Im}\delta)/D>0$. The transported $(0,1)$ graph coefficient is $\nu_\delta=\delta P/(2iD+\delta)$. The derivative at zero is $P/(2iD)$. Equivalently the smooth lift of $\partial_c$ in moving coordinates has vertical part $r_1\partial_{w_2}$ and $\bar\partial r_1=-(d\bar\zeta_2-a d\bar\zeta_1)/(2iD)$, giving the stated sign.

A nonzero constant tangent-valued $(0,1)$ form on a compact torus cannot be $\bar\partial$ of a global vector field. Write that field in the invariant holomorphic frame and integrate each derivative of its periodic coefficients against invariant probability volume. Every such derivative has integral zero, whereas the coefficient of $\partial_{\zeta_2}\otimes d\bar\zeta_2$ above has nonzero integral. This proves nonvanishing for the fibre deformation.

We explain its transfer to the threefold. Over the dual numbers $A=\mathbb C[\epsilon]/(\epsilon^2)$ at $c$, the chart proof of \ref{fam:src:A.1.1} gives the relative line isomorphism $\omega^{-2}_{X_A/A}=f_A^*\mathcal O_B(p_2)$. Extra parameter differentials disappear in relative forms and all determinant characters are independent of $c$. The two pulled-back sections form an $A$-basis of $H^0(X_A,\omega^{-2})$: reduce a section modulo $\epsilon$, subtract its linear combination of the two lifts, and use the exact sequence $0\to L\xrightarrow{\epsilon}L_A\to L\to0$ to write the remainder uniquely as $\epsilon$ times another such combination. The same argument proves independence.

A first-order trivialization would preserve this complete system and give a base automorphism reducing to the identity. It fixes all three special values even over $A$. On the finite etale charts the relative critical ideal is $(s_j^{m_j-1})$, and $t_j=s_j^{m_j}$ is zero on that scheme; its constant parameter ring $A$ injects. At a cusp triple chart, the critical ideal of $t_c=z_0z_1z_2$ is $(z_0z_1,z_0z_2,z_1z_2)$ and has the same property. There are no other critical values. Therefore an infinitesimal displacement of a special value would remain nonzero in its critical chart, a contradiction. Writing a base infinitesimal automorphism as $t+\epsilon(b_0+b_1t+b_2t^2)$, fixing $0,1,\infty$ gives $b_0=b_1=b_2=0$. The trivialization would consequently trivialize the displayed nonzero torus deformation. Finally a zero Kodaira--Spencer class produces such a first-order trivialization by subtracting from local product charts the vector fields whose coboundary is the differentiated transition cocycle. This proves the claim.
\end{proof}

\begin{proposition}\label{fam:src:A.4.3}\label{fam:st:A.4.3}
The non-isotriviality argument of Viaclovsky, \cite[Proposition 4.7]{ViaclovskyLog2026}, follows for the present family from \ref{fam:src:A.4.2} and \ref{fam:src:A.1.1}. The cited countability argument proves local non-isotriviality and positive Kuranishi dimension; the computation in \ref{fam:src:A.4.2} additionally proves nonzero derivative at every admissible parameter. The arithmetic method of Yang, \cite[Lemmas 6.6--6.9 and Proposition 6.15]{Yang2026}, specializes to the same centralizer, finite-norm and collar-gluing operations after replacing his displayed matrices by $A_1,A_2$; those replacements are evaluated in \ref{fam:src:A.2.1} and \ref{fam:src:A.3.1}.
\end{proposition}
\Uses{the hypotheses, notation and definitions stated here; \cite[Proposition 4.7]{ViaclovskyLog2026}; \cref{fam:src:A.4.2}; \cref{fam:src:A.1.1}; \cite[Lemmas 6.6--6.9 and Proposition 6.15]{Yang2026}; \cref{fam:src:A.2.1}; \cref{fam:src:A.3.1}}

\begin{proof}
On a fixed ordinary fibre, the three last periods remain fixed and span $\mathbb C^2$, whereas the first changes by $\delta e$. A marked complex-linear identification fixing those three periods must be the identity, so the marked period map is nonconstant. The set of integral markings of one torus is countable. A holomorphic nonconstant map from a disc has uncountable image; hence the unmarked fibre family is locally non-isotrivial. \ref{fam:src:A.1.1} transfers this to total spaces. This is exactly the implication used in the cited Proposition 4.7, whose statement and proof we have read. A nonconstant classifying map gives positive dimension of its local target. It does not by itself determine the derivative at the base point, as the explicit holomorphic map $z\mapsto z^2$ has nonconstant image and derivative zero at zero. \ref{fam:src:A.4.2} computes the derivative here.

Yang's cited lemmas use an integral commuting lattice map, an affine fibre lift, normalization of the root base change, and a norm-lattice defect. We have exhibited the respective maps and norm equations for the present matrices. His Proposition 6.15 obtains $6/r$ from the first-column scaling and its order-six norm image. Here the first-column scaling is one and the two computed norm images impose parity, yielding two. The common proof operations transfer through these explicitly stated matrix substitutions; no conclusion about one family is inferred merely from the label of the other. His nonzero absolute Kodaira--Spencer conclusion is obtained in the present setting by the explicit computation and intrinsic-pencil transfer above.
\end{proof}

\subsection{The assembled theorem}\label{fam:section:A.5}

\begin{theorem}\label{fam:src:A.5.1}\label{fam:st:A.5.1}\label{wb:pshift:thm:classification}
Fix the original cusp-regular period solution, the clockwise peripheral marking, the prescribed periodic triangular cusp fan, the affine vectors $(\varepsilon,-\varepsilon')$, and the untwisted cusp collar. For every $c,c'\in\mathcal P$, the constructed compact complex threefolds satisfy
\[
 X_c\simeq X_{c'}\quad\Longleftrightarrow\quad c'-c\in2\mathbb Z.
\]
Every biholomorphism is over the identity of the original $\mathbb P^1$. The proper holomorphic family of their isomorphism classes has parameter $0<|w|<e^{-\pi M}$ and nonzero Kodaira--Spencer class at every point.
\end{theorem}
\Uses{the hypotheses, notation and definitions stated here; \citet[Theorem~1.4]{DuEtAl2026v1}; \cref{app:rederive:fam-src-A.5.1}}


The standard classification is \citet[Theorem~1.4]{DuEtAl2026v1}. A complete re-derivation is given in Appendix~\ref{app:rederive:fam-src-A.5.1}.


\begin{theorem}\label{fam:src:A.5.2}\label{fam:st:A.5.2}\label{inv:D1.3}\label{inv:D2.1}\label{inv:D3.1}\label{cdp:st:E9}
For every admissible $c$, the tangent cohomology of the actual compact
member is
\[
 \bigl(h^0(X_c,T_{X_c}),h^1(X_c,T_{X_c}),h^2(X_c,T_{X_c}),
 h^3(X_c,T_{X_c})\bigr)=(1,1,1,0).
\]
Its Kuranishi germ is a smooth reduced disc. The family of \ref{fam:src:A.5.1} is
miniversal at every parameter: its Kodaira--Spencer map is an isomorphism,
and every sufficiently small deformation is induced from this family.
More precisely, for the intrinsic fibration $f:X_c\to\mathbb P^1$,
\[
 f_*T_{X_c}=\mathcal O\oplus\mathcal O(-1),\qquad
 R^1f_*T_{X_c}=\mathcal O\oplus\mathcal O(-1)^2,\qquad
 R^2f_*T_{X_c}=\mathcal O.\tag{D.4}
\]
The Leray class $\xi\otimes[\gamma]$ generates $H^1(X_c,T_{X_c})$.
The primary Kuranishi bracket is zero; the one-dimensional obstruction
space $H^2(X_c,T_{X_c})$ therefore does not imply a nonzero obstruction
map.

\end{theorem}
\Uses{the hypotheses, notation and definitions stated here; \cref{fam:src:A.5.1}; \citet[Theorem~6.12 and Proposition~6.16]{DuEtAl2026v1}; \cref{app:rederive:fam-src-A.5.2}; \cref{fam:section:B.2}; \cref{fam:section:B.3}; \cite[Sections 2--3 and Proposition 6.1]{ViaclovskyTwo2026}; \cref{fam:eq:ab-normal-quotient}}


The tangent dimensions and Kuranishi germ are re-derived from \citet[Theorem~6.12 and Proposition~6.16]{DuEtAl2026v1}. A complete re-derivation is given in Appendix~\ref{app:rederive:fam-src-A.5.2}.


\paragraph{Scope and a precise changed-input test.}
The hypotheses of \ref{fam:src:A.5.1}--\ref{fam:src:A.5.2} fix the original period system and
multiplicities $(3,4)$; \ref{fam:section:B.2} treats its arbitrary affine and cusp data,
and \ref{fam:section:B.3} classifies its stipulated higher-multiplicity variants.
A concrete extension beyond these hypotheses begins with
$y^2=x^3+tx+1$ in \cite[Sections 2--3 and Proposition 6.1]{ViaclovskyTwo2026}.
Already one rank-one filling over a specified root $p_1$ of
$4t^3+27=0$ replaces the hexagonal normalization calculation by
\[
 \begin{gathered}
 E'=\mathbb C^*/\langle c_0\rangle,\qquad
 R=\mathbb P_{E'}(\mathcal O_{E'}\oplus L),\\
 R=(\mathbb P^1\times\mathbb C^*)/
       \langle(z,v)\mapsto(\gamma_0z,c_0v)\rangle.
 \end{gathered}
\]
with its boundary sections identified by $v\mapsto\beta_0v$.
Here $c_0,\gamma_0,\beta_0$ are exactly the nonzero central multipliers
in that paper's equation (3.1). The invariant nodal volume form there
trivializes the ambient canonical line, so the branch-normal map has
quotient $\nu_*K_R$ by the same local coefficient map as
\eqref{fam:eq:ab-normal-quotient}. The ruling has
$\pi_*\mathcal O_R=\mathcal O_{E'}$ and $R^1\pi_*\mathcal O_R=0$;
Leray and Serre duality therefore give $h^\bullet(K_R)=(0,1,1)$.
The first unevaluated global datum for transferring the deformation
argument is the connecting map from this single
$H^1(R,K_R)=\mathbb C$ summand into $H^2(T_X(-\log D))$, where
$D$ is the union of the three reduced singular fibres of that model.
Its direct attack is to compute the boundary-translation Cech cocycle
using the full holomorphic unit functions in that equation (3.1),
their central values $(\gamma_0,c_0,\beta_0)$, and the corresponding
vertical frames.
The global identification argument also requires the integral
centralizer and collar defects of that model's four monodromy matrices.
These are specified substitutions to perform for the stated one-fibre
test; no classification or deformation conclusion for that changed
period system is included in \ref{fam:src:A.5.1}--\ref{fam:src:A.5.2}.

\section{Discrete data and the geometric parameter}\label{fam:section:B}

\subsection{All affine vectors and cusp sections}\label{fam:section:B.2}

\begin{proposition}\label{fam:src:B.2.1}\label{fam:st:B.2.1}
Let
\[
 w_1=r_1\varepsilon+k_1d,\qquad w_2=r_2\varepsilon'+k_2d,
 \qquad 3\nmid r_1,\quad 2\nmid r_2,
\]
and use an arbitrary holomorphic cusp collar section $\sigma$ in the construction of \cite[Section 6.4]{AC2026}. Write $l_0=\gamma(c(\sigma))$ in its positive cusp-winding convention, and set $p=12l_0-4r_1-3r_2$. Let primes denote a second such datum with the same period solution and cusp fan. Then
\[
 X(c;w_1,w_2,\sigma)\simeq X(c';w'_1,w'_2,\sigma')
\]
holds exactly for the following arithmetic conditions:
\[
 n=c'-c\in\mathbb Z,\qquad
 p'=a p\text{ for some }a\in\{1,-1\},\qquad
 n\equiv k_2-k'_2\pmod2.
\]
All maps are over $B$. A map with the indicated $a$ has fibre-linear part $aI_2$. In particular the complete invariants are
\[
 \bigl(|p|,\ e^{-\pi i(c+k_2)}\bigr),\qquad
 |p|\ge1,\quad\gcd(|p|,6)=1.
\]
Every value of this positive integer occurs. The full collar section enters the construction, but its additional analytic part and its second integral winding contribute no further invariant.
\end{proposition}
\Uses{the hypotheses, notation and definitions stated here; \cite[Section 6.4]{AC2026}; \cref{fam:src:A.1.1}; \cref{fam:src:A.3.1}}

\begin{proof}
The determinant calculation in \ref{fam:src:A.1.1} was independent of translations, so the base is still intrinsic. The centralizer proof gives $c'=c+n$ and $aI_2$ exactly as before. Put $S_n=I+nE$ and compare all target periods in source coordinates. On the finite root cover the exact deck equation is
\[
 B_j(g_jz)-R_jB_j(z)=\Pi_c(g_jz)\left(\frac{d_j}{m_j}+\ell_j\right),
 \quad d_j=S_nw'_j-a w_j,
 \quad N_j\ell_j=-d_j.                                      \tag{B2.1}
\]
Necessity follows by the normalized-base-change lift already proved. Conversely the norm equation has a rational solution $r_j^*$ of
$(I-A_j)r_j^*=d_j/m_j+\ell_j$, since $\ker N_j$ over $\mathbb Q$ equals the image of $I-A_j$. We can choose $\gamma(r_j^*)=0$: the columns $\widehat u,\widehat w$ already span that image over $\mathbb Q$. Then $B_j=\Pi_cr_j^*$ is holomorphic on the whole root disc and solves the exact equation. Thus the norm condition is sufficient for a local isomorphism, not only for conjugacy on the central real torus.

The computed norm images translate (B2.1) into
\[
 r'_1\equiv ar_1\pmod3,\quad r'_2\equiv ar_2\pmod4,
 \quad ak_2-k'_2-nr'_2\equiv0\pmod2.                       \tag{B2.2}
\]
There is no condition on $k_1$. Since $r'_2$ is odd, the last condition is the displayed parity. On the finite collar the middle translation is
$A_j=B_j+(\log s_j)/(2\pi i)\,\Pi_cd_j$ and has integral defect $\ell_j$. Its $\gamma$ defect is
$\alpha_j=(ar_j-r'_j)/m_j$.

We prove the precise cusp necessity used next. A holomorphic isomorphism between cusp neighbourhoods with fibre-linear part $aI_2$ has, after lifting near a generic smooth central point to the toric covers, the form
\[
 (x,t_c)\longmapsto(t_c^q u(t_c)x^a,t_c),\qquad q\in\mathbb Z^2,
 \quad u(0)\in(\mathbb C^*)^2.                           \tag{B2.3}
\]
Here the assertion concerns this local lift near a generic central point. To prove it, take a source chart near the generic point of its component and a target chart near its image. The etale toric quotient permits a local lift over a simply connected neighbourhood. The torus characters in a chart at the height-one ray $v$ are $x_i=t_c^{v_i}X_i$, with $X_i$ invertible. On $t_c\ne0$ the affine fibre equation gives $x'_i/x_i^a$ independent of the $X$ coordinates. Its expression in the two charts is therefore $t_c^{v'_i-av_i}$ times a holomorphic unit of $t_c$, including zero. This proves (B2.3). Taking logarithms shows that its winding is in $\Lambda_{\rm tor}=\mathbb Z\widehat w+\mathbb Zd$. Changing a lift adds $(I-A_0)\Lambda=\Lambda_{\rm tor}$, so the vanishing of its class in $\bar\Lambda=\Lambda/\Lambda_{\rm tor}$ is intrinsic. This proof supplies the necessary extension test; no assertion about arbitrary punctured translations is used.

We retain the actual collar convention. The map from the cusp piece to the middle family is $T_{-\sigma}G_0^{-1}$. Consequently the middle translation induced by (B2.3) is
\[
 A_0=sq+a\sigma-\sigma'+\eta_0,
\]
with $\eta_0$ holomorphic on the cusp disc after a logarithm of the unit. For the clockwise generator $s\mapsto s-1$, the class of its defect in $\bar\Lambda$ is $-a c(\sigma)+c(\sigma')$. On the other hand every middle integral cocycle satisfies
$\ell_0=-A_0(\ell_1+A_1\ell_2)$, whose $\gamma$ coordinate is $-\alpha_1-\alpha_2$. Equality of these coordinates is exactly
\[
 l'_0-a l_0=-(ar_1-r'_1)/3-(ar_2-r'_2)/4,
\]
or $p'=ap$. This proves the remaining necessity. Notice that $p\equiv-r_1\pmod3$ and $p\equiv r_2\pmod2$, so $p$ is nonzero and coprime to six.

For sufficiency, $p'=ap$ implies the first two congruences of (B2.2) by reduction modulo three and four. The parity supplies the remaining one. Choose integral solutions $\ell_1,\ell_2$ of (B2.1). Their $\gamma$ coordinates already have the required values. Adding any element of $\ker N_1=(I-A_1)\Lambda$ preserves (B2.1). In particular adding an arbitrary integer multiple of $e_1=(0,1,-2,0)^t$ changes the $\widehat u$ coordinate of $\ell_0$ by an arbitrary integer while leaving its $\gamma$ coordinate fixed. We can therefore arrange
$[\ell_0]=-a c(\sigma)+c(\sigma')$ in the entire rank-two group $\bar\Lambda$. Choose lifts of $\sigma,\sigma'$ on the cusp half-plane. Their remaining defect differs from $\ell_0$ by an actual toric lattice vector; an integral $q$ makes $sq+a\sigma-\sigma'$ have exactly defect $\ell_0$.

The finite translations and this cusp reference now define a locally nonempty affine torsor with linear difference sheaf precisely $\mathcal V$ of \ref{fam:src:A.3.1}. The terms $a\sigma-\sigma'$ are retained, even when their chosen lifts have essential singularities at the puncture; they belong to the specified collar reference, while differences of allowed extensions are holomorphic $\eta_0$. The proof of $H^1(B,\mathcal V)=0$ gives a global middle translation. The finite maps descend by (B2.1), and the cusp map is $x\mapsto t_c^qu(t_c)x^a$. The fan is preserved by all integer translations and by $y\mapsto-y$; exponentials of the integer period shift agree. Hence this cusp map and its inverse extend over every stratum. The finite overlap is the equality defining $A_j$, and the cusp overlap is the equality defining $A_0$. Thus the maps glue globally.

Finally parity is equivalent to equality of the two displayed exponentials. Passing from $p$ to $-p$ is effected by $a=-1$ and the same proof; for fixed sign the congruences have just been shown sufficient. Every positive integer coprime to six is obtained by choosing $r_1\in\{1,2\}$ and odd $r_2$ in the appropriate residue classes modulo three and four, and then taking $l_0=(p+4r_1+3r_2)/12$. This also proves that the extra collar data have been accounted for, rather than omitted from the statement.
\end{proof}

\begin{proposition}\label{fam:src:B.2.2}\label{fam:st:B.2.2}
For multiplicities $(3,4)$ and $|p|=1$, every admissible discrete member is biholomorphic to one member of \ref{fam:src:A.5.1}, uniquely modulo $2\mathbb Z$. At fixed nonzero admissible Engel parameter $u$ there are exactly two isomorphism classes, exchanged by a positive circuit of $u$ about zero. In the convention that Engel's root generator is the inverse of our clockwise one and his circle class is $d$, the identity of the middle quotient sends
\[
 X_{c+1}\longleftrightarrow X(u;a_0-d/3,\ a_1+d/4).
\]
Reversing both declared root generators reverses these two displayed signs.
\end{proposition}
\Uses{the hypotheses, notation and definitions stated here; \cref{fam:src:A.5.1}; \cref{fam:src:B.2.1}; \cref{fam:src:A.2.1}}

\begin{proof}
Take the canonical target data $(r'_1,r'_2,k'_1,k'_2,l'_0)=(1,-1,0,0,0)$, for which $p'=-1$. Choose $a$ so that $p'=ap$, and an integer $n\equiv k_2\pmod2$. \ref{fam:src:B.2.1} supplies the isomorphism to $X_{c+n}$. The uniqueness is \ref{fam:src:A.5.1}. At fixed $u=e^{-2\pi ic}$, its logarithms differ by integers; the invariant $e^{-\pi i(c+k_2)}$ takes exactly the two square roots of $u$. Both occur, already as $X_c,X_{c+1}$, and \ref{fam:src:A.2.1} distinguishes them.

We spell out the integral comparison to Engel's common basis $B_0=(e_1,e_2,e_3,e_4)$, retaining his $C f_4=e_1+e_4$:
\[
 P_{\rm E}=\begin{pmatrix}0&-2&-1&-2\\2&-1&0&-1\\0&2&1&3\\-1&0&0&0\end{pmatrix},\qquad
 T_0P_{\rm E}=P_{\rm E}A_1,\quad T_1P_{\rm E}=P_{\rm E}A_2.
\]
Its determinant is $-1$, $P_{\rm E}d=\delta_0$, $\psi P_{\rm E}=-\gamma$, $P_{\rm E}\varepsilon=-e_4$ and $P_{\rm E}\varepsilon'=\delta_0-Cf_4$. These are exact lattice homomorphism identities, verified independently through Viaclovsky's intermediate basis in \texttt{ab\_geometry.py}. They also agree with the complex periods: at $z_1$ the complex-linear map
$(\zeta_1,\zeta_2)\mapsto(-\zeta_1/(\rho+1),\zeta_2+\zeta_1/3)$ carries $\Pi_cP_{\rm E}^{-1}$ to Engel's columns $(1,0),(\rho-1,0),((1+\rho)/3,1/3),(0,-\kappa_1)$. At $z_2$ his $Cf_4$ has preimage $d-\varepsilon'$ and fixed period $1-\kappa_2$. Thus the comparison fixes the positive circle class and the actual local complex structures, in addition to their monodromy matrices.

For the sign assertion, inverting the concrete clockwise affine generator changes its invariant translation $\Pi_cv_j/m_j$ to $-\Pi_cv_j/m_j$. Thus our dictionary is $a_{j-1}=-P_{\rm E}v_j/m_j$ in Engel's lattice marking (and its period image in the torus). Integer shift gives $\Pi_{c+1}v_j=\Pi_c(v_j+\gamma(v_j)d)$, with $\gamma(v_1)=1,\gamma(v_2)=-1$, proving the two signs. A positive $u$ circuit changes $\log u$ by $2\pi i$ and hence changes $c$ by $-1$, exchanging the square roots. This gives the claimed permutation with its orientation.
\end{proof}

\begin{proposition}[counterexample, with corrected family]\label{fam:src:B.2.3}\label{fam:st:B.2.3}
The assertion in \cite[Remark 8.1]{Engel2026} of a global proper family over the original punctured $u$-disc for fixed sphere logarithmic data requires a double-cover qualification. The retained family with such data does not descend through $u=w^2$. It does give a proper holomorphic family on each simply connected $u$-open set, and a global proper family over the punctured $w$-disc as in \ref{fam:src:A.3.1}.
\end{proposition}
\Uses{the hypotheses, notation and definitions stated here; \cite[Remark 8.1]{Engel2026}; \cref{fam:src:A.3.1}; \cref{fam:src:A.2.1}; \cref{fam:src:B.2.1}; \cref{fam:src:A.4.1}}

\begin{proof}
Descent through $u=w^2$ would furnish an isomorphism between the two pulled-back fibres at $w$ and $-w$, namely $X_c$ and $X_{c+1}$. \ref{fam:src:A.2.1} excludes that isomorphism by the nonzero class of $d$ in $\Lambda^{A_2}/N_2\Lambda$. This is an exact order-two obstruction on the normalized finite cover. It applies to all sphere data by \ref{fam:src:B.2.1} because $r_2$ is odd.

There is also a direct analytic witness near zero. The relative Albanese invariant of the multiplicity-four reduced fibre in a holomorphic descended family would be a single-valued holomorphic function of $u$. After pullback it would be even in $w$. But \ref{fam:src:A.4.1} gives $j_2(w)=C^{-1}w^{-1}+744+O(w)$, where $C=e^{-\pi i\kappa_2^{\rm part}}\ne0$; its Laurent coefficient of $w^{-1}$ is nonzero. Thus it is not even on any punctured neighbourhood. This verifies the obstruction independently of a choice of isomorphism lift. Local inverse logarithms of $u$ give the proper local families already constructed; coherent even shifts give their global double-cover family. The untwisted middle quotient itself has the integer shear identification, so this refutation concerns the fixed sphere logarithmic data in the stated remark.
\end{proof}

\subsection{Higher multiplicities}\label{fam:section:B.3}

\begin{proposition}\label{fam:src:B.3.1}\label{fam:st:B.3.1}
Let $k,l$ be positive integers with $\gcd(3k,4l)=1$, and put $m_1=3k,m_2=4l$. Consider Viaclovsky's sphere members expressed in the same clockwise marking, with $w_j=r_j\varepsilon_j+s_jd$, $\gcd(r_j,m_j)=1$, cusp winding $l_0$, and $p=m_1m_2l_0-m_2r_1-m_1r_2=\pm1$, where $\varepsilon_1=\varepsilon,\varepsilon_2=\varepsilon'$. Two such members are biholomorphic precisely when their multiplicity pair agrees and, for some $a\in\{1,-1\}$ and integer $n=c'-c$,
\[
 p'=ap,\quad r'_j\equiv ar_j\pmod{m_j},\quad
 nr'_1\equiv as_1-s'_1\pmod{k},\quad
 nr'_2\equiv as_2-s'_2\pmod{2l}.
\]
The congruences on $r'_j$ already follow from $p'=ap$. For fixed discrete vectors and cusp collar the identification subgroup is $2kl\mathbb Z$. For fixed admissible $u$ and this multiplicity pair, all the stated sphere translation data give exactly $2kl$ isomorphism classes, permuted in one cycle by a circuit of $u$ about zero.
\end{proposition}
\Uses{the hypotheses, notation and definitions stated here; \cref{fam:src:A.1.2}; \cref{fam:src:A.2.1}; \cref{fam:src:B.2.1}}

\begin{proof}
We first recover the base for these sphere members without using the degree-two canonical pencil. A biholomorphism $F:X\to X'$ gives $g=f'F$. Since $H^2(S^6,\mathbb R)=0$, the closed form $g^*\omega_{\rm FS}$ is exact. Restrict to a regular torus $T$ and wedge with a flat Kahler form on $T$. Stokes' theorem gives integral zero, while the integrand is nonnegative and strictly positive wherever $dg|_T\ne0$. Thus $g$ is constant on each regular fibre. The image of $(f,g)$ is a compact analytic curve whose projection to $B$ has degree one; its normalization gives a holomorphic $h:B\to B$ with $g=hf$. Apply the same construction to $F^{-1}$ to see that $h$ is an automorphism. The three homological local monodromy orders are three, four and infinity, independently of $k,l$; their intertwining maps under $F$ therefore fix each special point. Hence $h=1$. The nonreduced fibre multiplicities now force equal $k,l$. This is the full reduction in the read Proposition 4.7 of Viaclovsky's paper, written here with its exact maps.

The root model of degree $m_1$ is the pullback of the degree-three model by $s\mapsto s^k$, and likewise at the other point. Its linear action remains $A_1$ or $A_2$, while its root generator has order $m_j$. The primitive invariant coordinate $r_j/m_j$ proves freeness of every nonidentity power. The normalization and deck-commutation arguments of \ref{fam:src:A.1.2} and \ref{fam:src:A.2.1} apply to these root covers without changing a generator. Their norms are
\[
 N_{3k}=kN_1,\qquad N_{4l}=lN_2,
\]
so the fixed-coordinate norm images are $3k\mathbb Z\varepsilon+k\mathbb Zd$ and $4l\mathbb Z\varepsilon'+2l\mathbb Zd$. The norm equation (B2.1), with $m_j$ replaced by these actual covering degrees, gives exactly the displayed congruences. The centralizer and first-column comparison remain unchanged.

For completeness the integer $p$ records the primitive sphere relation in this marking. The common torus coinvariants are detected by $\gamma$, since $(A_1-I)\Lambda+(A_2-I)\Lambda=\ker\gamma$ by their explicit saturated bases. The finite meridian relations are $x^{m_1}=c^{r_1}$, $y^{m_2}=c^{r_2}$, and the cusp relation is $xy=c^{l_0}$ with $c$ central. Eliminating $y$ gives an abelian relation matrix with determinant $-p$ and primitive first elementary divisor because $\gcd(m_1,m_2)=1$. Thus the sphere data have $p=\pm1$. Reducing $p'=ap$ modulo $m_1,m_2$ gives the two congruences for $r_j$.

To prove sufficiency, use the local rational solutions of the norm equation and the identical cusp comparison in \ref{fam:src:B.2.1}. The first-coordinate extension on the new root discs has vanishing orders $2k$ and $l$, precisely the orders of the pulled-back $\Theta$. Dividing by it again gives $\mathcal O_B(-p_0)$, and finite averaging again gives the same rank-two extension sheaf. Thus $H^1(B,\mathcal V)=0$ still removes the analytic cocycle. The images of $I-A_j$ and their arbitrary $\widehat u$ adjustments are unchanged, so $p'=ap$ is exactly the remaining cusp winding condition. Every local map and overlap is the one already written in \ref{fam:src:B.2.1}, with the new denominator $m_j$; this proves global sufficiency.

For identical data, $p\ne0$ excludes $a=-1$. The congruences reduce to $k\mid nr_1$ and $2l\mid nr_2$. The two primitivity assumptions imply $k\mid n$ and $2l\mid n$. Since $\gcd(k,2l)=1$, their least positive common multiple is $2kl$. For different sphere translation data, $r'_1$ is invertible modulo $k$, $r'_2$ modulo $2l$, and the Chinese remainder theorem gives exactly one integer shift class modulo $2kl$. The values of $s_1$ modulo $k$ and $s_2$ modulo $2l$ all occur without affecting freeness or $p$. Thus there are exactly $2kl$ classes at fixed $u$. The change $c\mapsto c-1$ on a positive $u$ circuit acts by a generator of this cyclic residue group. Program \texttt{ab\_exact.py} verifies the repeated norms and shift steps for five coprime pairs; the preceding symbolic identities prove all pairs.
\end{proof}

\subsection{Automorphisms and the remaining real-structure calculation}\label{fam:section:B.4}

\begin{proposition}\label{fam:src:B.4.1}\label{fam:st:B.4.1}
For each member of the retained $(3,4)$ family, $\operatorname{Aut}(X_c)$, including all its components, is $\mathbb C/\mathbb Z\simeq\mathbb C^*$. Its action on the ordinary fibres is $\zeta\mapsto\zeta+be$, and the exponential group isomorphism is $b\mapsto e^{2\pi ib}$.
\end{proposition}
\Uses{the hypotheses, notation and definitions stated here; \citet[Theorem~1.4]{DuEtAl2026v1}; \cref{app:rederive:fam-src-B.4.1}}


The full automorphism group is \citet[Theorem~1.4]{DuEtAl2026v1}. A complete re-derivation is given in Appendix~\ref{app:rederive:fam-src-B.4.1}.



\begin{proposition}\label{fam:src:B.4.2}\label{fam:st:B.4.2}
Use the normalization $(\beta^{\rm part}+\tau)(p_0)=0$. There is an antiholomorphic diffeomorphism $X_c\to X_{c'}$ exactly when
\[
 c'-(1-\bar c)\in2\mathbb Z,
 \qquad\text{equivalently }w'=-\bar w.
\]
Consequently $\overline{X_c}\simeq X_{1-\bar c}$. The members admitting antiholomorphic self-maps are exactly those with $\operatorname{Re}c\in\tfrac12+\mathbb Z$, equivalently $w=-\bar w$. On each such member the group of holomorphic and antiholomorphic diffeomorphisms is
\[
 \operatorname{Aut}^{\pm}(X_c)\simeq\mathbb C^*\rtimes\mathbb Z/2,
 \qquad z\longmapsto1/\bar z
\]
for the action of the second factor. Its antiholomorphic maps form the nonidentity coset. For all other members the full group is the holomorphic group of \ref{fam:src:B.4.1}.
\end{proposition}
\Uses{the hypotheses, notation and definitions stated here; \cref{fam:src:B.4.1}; \cref{fam:src:A.3.1}; \cref{fam:src:B.2.1}; \cref{fam:src:A.5.1}}

\begin{proof}
We first construct the comparison before returning its affine data to the canonical choice. Lift conjugation of the real orbifold base to the antiholomorphic involution $\iota$ of its universal cover whose cusp branch has $\tau(\iota z)=-\overline{\tau(z)}$. A lift exists by the universal covering property; it extends at the finite orbifold points, whose local equations are $s_j^{m_j}=t_j$. Choose the usual paths from a real point $t>1$, so reflection fixes the second centre and sends the first chosen lift to $g_2^{-1}z_1$. The induced involution of the deck group is
\[
 \phi(g_1)=g_2^{-1}g_1^{-1}g_2,\qquad
 \phi(g_2)=g_2^{-1},\qquad \phi(g_0)=g_0^{-1}.
\]
These identities also follow by reflecting the lifted peripheral paths; their product and their squares agree with $g_1g_2g_0=1$. The cusp branch fixes the lift uniquely.

All reflection identities for the periods can be checked before invoking uniqueness. Denote their affine transformation formulas by $G_1,G_2$ on the triple $(\tau,\mu,\beta)$ and put $K(\tau,\mu,\beta)=(-\bar\tau,1-\bar\mu,-\bar\beta)$. Direct rational substitution gives
\[
 KG_1K=G_2^{-1}G_1^{-1}G_2,\qquad KG_2K=G_2^{-1}.
\]
For example $G_2^{-1}(\tau,\mu,\beta)=(-1/\tau,(1-\mu)/\tau,\beta+3-6(1-\mu)^2/\tau)$ and
$G_1^{-1}(\tau,\mu,\beta)=(1/(1-\tau),1-\mu/(1-\tau),\beta-2+6\mu^2/(1-\tau))$. These formulas verify every constant in the two equalities. The reflected $\mu$ is cusp bounded and has the same affine laws, so their difference is a global section of $\mathcal O_B(-p_0)$ and vanishes. After this equality, the difference between the reflected and original $\beta^{\rm part}$ is a holomorphic invariant base function, hence constant; its cusp value is zero under the stated normalization. Therefore
\[
 \mu(\iota z)=1-\bar\mu(z),\qquad
 \beta^{\rm part}(\iota z)=-\overline{\beta^{\rm part}(z)}.
\]
The original $\tau$ identity follows from its $j$ equation and its specified cusp branch. In particular $D_{-\bar c}(\iota z)=D_c(z)$, so this comparison stays in the exact admissible domain.

Let
\[
 H=\begin{pmatrix}1&0&0&0\\0&1&0&0\\-6&0&-1&0\\0&-1&0&-1\end{pmatrix}.
\]
Then $H^2=I$, $HA_2H=A_2^{-1}$, $HA_1H=A_2^{-1}A_1^{-1}A_2$, and
\[
 \Pi_{-\bar c}(\iota z)H=-\overline{\Pi_c(z)}.
\]
Thus $(z,\zeta)\mapsto(\iota z,-\bar\zeta)$ is a well-defined antiholomorphic torus map. It also descends under the middle deck action. Indeed covariance and the displayed lattice conjugacies imply $R_{\phi(g)}(\iota z)=\overline{R_g(z)}$ by multiplying the full real-invertible period matrices. This proves exact compatibility, not merely an equality of fibre lattices.

At the first finite centre compose this map with the target deck identification $g_2$ to return the centre $g_2^{-1}z_1$ to $z_1$. Its lattice map becomes $A_2H$ and it conjugates the clockwise generator to the inverse generator. At the second centre its lattice map is $H$ with the same reversal. Consequently the target affine vectors are exactly
\[
 w'_1=-A_2H\varepsilon=-\varepsilon+d,\qquad
 w'_2=-H(-\varepsilon')=\varepsilon'-3d.
\]
The minus signs arise because the inverse of a generator with invariant translation $w'_j/m_j$ has translation $-w'_j/m_j$. This calculation verifies the full affine actions on the root covers. For the overlaps choose the root coordinates of the real local base coordinates $t$ and $t-1$ compatible with these local reflections. This choice has an explicit comparison with any original root coordinate: their ratio is a holomorphic unit, so changing it changes the collar by $(2\pi i)^{-1}\log h\,\Pi w_j$, a translation extending on the entire root family. Its inverse gives a holomorphic isomorphism of the two glued descriptions. Thus the extension property and final classification transfer to the original coordinates. Then $\log s'_j=\overline{\log s_j}$ on paired branches. Under $\zeta\mapsto-\bar\zeta$, conjugation of $1/(2\pi i)$ supplies the compensating sign, and the identity of the two logarithmic collar translations follows exactly. Two roots of the same fixed real coordinate differ by an $m_j$th root of unity; this changes the comparison only by a constant period-coordinate translation, and the additive correction of \ref{fam:src:A.3.1} glues those constants. The preceding holomorphic-unit comparison handles a different original local base coordinate. Its inverse has the same construction.

At the cusp choose $s=(2\pi i)^{-1}\log t_c$, so $s\mapsto-\bar s$. The full holomorphic unit matrix satisfies
\[
 C_{-\bar c}(\bar t_c)=-\overline{C_c(t_c)}+
       \begin{pmatrix}6&0\\0&1\end{pmatrix}.
\]
The integral extra matrix has trivial exponential. Hence $(x,t_c)\mapsto(\bar x,\bar t_c)$ intertwines the full cusp deck groups, with the same integral index $\lambda$. The toric fan is defined over $\mathbb Z$, so this map extends antiholomorphically on every toric chart and stratum. It preserves the untwisted cusp collar. We have therefore constructed an antiholomorphic map to the actual target datum with parameter $-\bar c$, $r'_1=-1,k'_1=1,r'_2=1,k'_2=-3,l'_0=0$, whose $p'$ is $+1$.

Apply \ref{fam:src:B.2.1} to normalize this target back to the canonical $p=-1$ datum. Take $a=-1$ and the odd shift $n=1$. This gives a holomorphic map to $X_{1-\bar c}$. Every other antiholomorphic map, composed with the inverse of this one, is a holomorphic comparison between canonical members. \ref{fam:src:A.5.1} gives exactly the stated congruence. Exponentiating it gives $w'=-\bar w$.

It remains to verify the group extension rather than infer its square from $H^2$. For a real member, choose any resulting antiholomorphic self-map $J$. Its derivative on ordinary fibre coordinates is $+\text{complex conjugation}$: the initial reflection has derivative $-\text{complex conjugation}$, while the normalization just used has linear part $-I_2$. Hence it conjugates translation by $b\in\mathbb C/\mathbb Z$ to translation by $\bar b$. Its square is a holomorphic translation $d$. Commuting $J$ with $J^2$ shows $d=\bar d$ in $\mathbb C/\mathbb Z$. This equality forces the imaginary part of a representative of $d$ to vanish, so take $d\in\mathbb R$. Replacing $J$ by the composition of translation by $-d/2$ with $J$ changes its square to $d-d/2-d/2=0$. We have an antiholomorphic involution. Under $b\mapsto e^{2\pi ib}$, its conjugation action is $z\mapsto1/\bar z$. This proves the split group and specifies its antiholomorphic coset.
\end{proof}



\section{Classification and deformation questions}\label{push:section:E}
\subsection{Higher logarithmic transformations}\label{push:section:E.1}
We distinguish Viaclovsky's indices $0,1$ from the indices $1,2$ of
the question.  Put $m=3k$, $n=4l$, and assume $\gcd(m,n)=1$.
The lattice coordinates in this subsection are integral homology
coordinates, and matrices act on columns.

\begin{proposition}[the integral dictionary]\label{push:src:E.1.1}
Let $(\delta,a,b,d)$ be the marking in
\cite[Section 2, equations (2.1), (2.2), (2.7), (2.11)--(2.12)]{ViaclovskyLog2026},
and put $u_0=d+2a+2b$, $u_1=d+3a+3b$.
The homomorphism from the marking of the question to this marking is
\[
 P_V=\begin{pmatrix}
 -2&1&0&1\\-6&0&-1&0\\0&-1&0&0\\-1&0&0&0
 \end{pmatrix},\qquad \det P_V=-1.
\]
It intertwines the two finite monodromies and satisfies
\[
 P_V\widehat\delta=\delta,\qquad
 P_V\varepsilon=-u_0,\qquad
 P_V\varepsilon'=\delta-u_1.
\]
With the counterclockwise deck generators of that paper, its data
$m_0=3k$, $m_1=4l$ and
$a_i=[v_i/m_i]$ with $v_i=\rho_i\delta+N_i u_i$ correspond to
\[
 w_1=N_0\varepsilon-\rho_0\widehat\delta,
 \qquad
 w_2=N_1\varepsilon'-(\rho_1+N_1)\widehat\delta.
\]
For the section-preserving cusp comparison used there, $l_0=0$ and
$p=-(nN_0+mN_1)$.  Thus its sphere condition is
$b:=nN_0+mN_1=\pm1$.
The brackets in $a_i$ mean the torsion point in the smooth
root torus, using its section-preserving origin and the
transported period vector.  Replacing $v_i$ by an integral
multiple $m_i$ times a period does not change that torsion
point, although its chosen lift fixes the collar bookkeeping.
\end{proposition}
\Uses{the hypotheses, notation and definitions stated here; \cite[Section 2, equations (2.1), (2.2), (2.7), (2.11)--(2.12)]{ViaclovskyLog2026}}

\begin{proof}
In Engel's marking the four columns of the change to Viaclovsky's
basis are
\[
 Q_E=\begin{pmatrix}-2&1&0&-2\\-1&0&0&0\\3&-1&1&0\\0&0&0&1\end{pmatrix}.
\]
Direct integral multiplication gives $P_V=Q_E^{-1}P_E$ and
$V_0P_V=P_VA_1$, $V_1P_V=P_VA_2$, where $V_0,V_1$ are the
matrices in that paper's Lemma 2.1.  Multiplication by the three
specified invariant vectors gives the displayed identities.
Inverting a clockwise affine generator replaces its invariant
translation $w_j/m_j$ by $-w_j/m_j$.  Consequently the actual
counterclockwise translation vector is $-P_Vw_j$, which gives the
two formulas, including the otherwise easily missed $N_1$ in the
second coefficient.  Viaclovsky's reference section extends through
the smooth locus of the cusp; its comparison therefore has zero
class in the quotient lattice of toric directions.  This is precisely
$l_0=0$.  Substitution into the definition of $p$ proves its formula.
These matrix identities are independently executed in
\texttt{code/e\_push\_exact.py}, with their output attached.
\end{proof}

For a general cusp collar $\sigma$, write $c(\sigma)$ for
its residual translation defect in
$\overline\Lambda=\Lambda/(\mathbb Z\widehat w+
\mathbb Z\widehat\delta)$, using the convention
$\sigma(g_0z)-R_0\sigma(z)=\Pi_c(g_0z)\ell(\sigma)$
and $c(\sigma)=[\ell(\sigma)]$.
Set $l_0=\gamma(c(\sigma))$.  In the question's notation
$w_j=r_j\varepsilon_j+s_j\widehat\delta$, the integer is
$p=mn l_0-nr_1-mr_2$.
The section-preserving source data of \ref{push:src:E.1.1} have $c(\sigma)=0$.

\begin{theorem}[classification in Viaclovsky's coordinates]\label{push:src:E.1.2}
Use a local logarithm $c$ of the transported linearization,
$u=e^{-2\pi i c}$, and the same reference section and cusp comparison
on both sides.  Two sphere members with data
$(m,\rho_0,N_0;n,\rho_1,N_1;c)$ and primed data are biholomorphic
if and only if their multiplicities agree and there are
$\alpha\in\{1,-1\}$ and $h=c'-c\in\mathbb Z$ such that
\[
 b'=\alpha b,\qquad
 hN'_0\equiv\rho'_0-\alpha\rho_0\pmod{k},\qquad
 hN'_1\equiv\rho'_1-\alpha\rho_1\pmod{2l}.
\]
Every such biholomorphism covers the identity of $B$.
For fixed discrete data the period identification subgroup is
$2kl\mathbb Z$.  At fixed $u$ all sphere translation data give
exactly $2kl$ classes.  A positive $u$-circuit acts on them by a
single cycle.
\end{theorem}
\Uses{the hypotheses, notation and definitions stated here; \cref{push:src:E.1.1}}

\begin{proof}
We first give the complete sufficiency argument in the stated problem's
clockwise coordinates, writing $w_j=r_j\varepsilon_j+s_j\widehat\delta$.
Set $S_h=I+hE$ and $d_j=S_hw'_j-\alpha w_j$.
The repeated norm maps have images
\[
 N_m\Lambda=m\mathbb Z\varepsilon+k\mathbb Z\widehat\delta,
 \qquad
 N_n\Lambda=n\mathbb Z\varepsilon'+2l\mathbb Z\widehat\delta.
\]
This follows directly from
\[
 N_3v=3\gamma(v)\varepsilon+(2v_u+v_w+3v_\delta)\widehat\delta,
 \quad
 N_4v=4\gamma(v)\varepsilon'+(2v_u+2v_w+4v_\delta)\widehat\delta.
\]
Suppose $p'=\alpha p$ and
$hr'_1\equiv\alpha s_1-s'_1\pmod k$,
$hr'_2\equiv\alpha s_2-s'_2\pmod{2l}$.
Reduction of $p'=\alpha p$ modulo $m,n$ gives
$r'_j\equiv\alpha r_j\pmod{m_j}$.
We may choose integral solutions of $N_{m_j}\ell_j=-d_j$ explicitly:
\[
 \begin{split}
 \ell_1&=\left(\frac{\alpha r_1-r'_1}{m},0,
 -\frac{s'_1+hr'_1-\alpha s_1}{k},0\right)^t,\\
 \ell_2&=\left(\frac{\alpha r_2-r'_2}{n},
 -\frac{s'_2+hr'_2-\alpha s_2}{2l},0,0\right)^t.
 \end{split}
\]
Put $y_j=d_j/m_j+\ell_j$; these vectors belong to the rational
kernel of $N_{m_j}$ and have first coordinate zero.
The following rational vectors solve $(I-A_j)b_j=y_j$:
\[
 \begin{split}
 b_1&=(0,(2y_{1u}+y_{1w})/3,(y_{1w}-y_{1u})/3,0)^t,\\
 b_2&=(0,(y_{2u}-y_{2w})/2,(y_{2u}+y_{2w})/2,0)^t.
 \end{split}
\]
The last coordinates of the equations follow from the displayed
norm-kernel equations.  The maps
\[
 (s,\zeta)\longmapsto(s,\alpha\zeta+\Pi_c(s)b_j)
\]
are holomorphic maps of the normalized torus families over the
degree-$m_j$ root discs.  Their inverses are holomorphic, and period
covariance and $(I-A_j)b_j=d_j/m_j+\ell_j$ give exact commutation
with the prescribed affine deck generators, modulo the integral
period $\ell_j$.  They therefore descend to biholomorphisms of the
finite fillings, including the full multiple fibres.

On the punctured root collar the corresponding middle translation is
\[
 A_j=\Pi_cb_j+\frac{\log s_j}{2\pi i}\Pi_cd_j.
\]
Its integral defect is $\ell_j$.  The middle cocycle relation fixes
$\ell_0=-A_0(\ell_1+A_1\ell_2)$.  In the quotient
$\overline\Lambda=\Lambda/(\mathbb Z\widehat w+\mathbb Z\widehat\delta)$,
the first-coordinate condition is exactly $p'=\alpha p$.
Adding an integral multiple of $(0,1,-2,0)^t\in\ker N_m$ to
$\ell_1$ changes the other quotient coordinate by a generator,
without changing this first coordinate.  We can consequently arrange
\[
 [\ell_0]=-\alpha c(\sigma)+c(\sigma')\quad\hbox{in }\overline\Lambda.
\]
An integral toric vector then supplies $q\in\mathbb Z^2$ such that
the cusp reference $sq+\alpha\sigma-\sigma'$ has precisely defect
$\ell_0$.  Here $s=(2\pi i)^{-1}\log t_c$ is the cusp logarithm;
the maps of cusp toric covers are
\[
 (x,t_c)\longmapsto(t_c^q a(t_c)x^\alpha,t_c),
 \qquad a(0)\in(\mathbb C^*)^2.
\]
Integer translations of the cocharacter lattice, its inversion,
and multiplication by a holomorphic unit preserve the specified
hexagonal fan and all its strata.  These maps and their inverses
extend holomorphically across the three double curves and both
triple points.

We justify the global analytic correction.  Differences of allowed
finite extensions are the equivariant holomorphic vertical vectors
on the complete root discs.  The extra subgroup of order $k$ at
the first point acts trivially on the vertical vector coefficients
and rotates $s$; its invariant holomorphic coefficients are exactly
those in $s^k$.  The remaining order-three equivariance is the
original one.  The argument at the second point uses $s^l$ and
order four.  Thus the coarse difference sheaf is literally the same
sheaf $\mathcal V_B$ as for multiplicities $(3,4)$, not an assumed
analogue.  Its first-coordinate quotient is
$\mathcal O_B(-p_0)$: the modular frame has vanishing orders $2k,l$
on the two enlarged roots and a simple cusp pole.  Its fixed
second-coordinate subline is $\mathcal O_B$.  Finite averaging
provides local lifts, hence
\[
 0\longrightarrow\mathcal O_B\longrightarrow\mathcal V_B
 \longrightarrow\mathcal O_B(-p_0)\longrightarrow0.
\]
Both end terms have vanishing $H^1$ on $\mathbb P^1$, so
$H^1(B,\mathcal V_B)=0$.  The permitted local translations form a
torsor under this sheaf: on each overlap their defects agree, so
their difference is a single-valued section of the indicated
difference sheaf.  Vanishing of its \v Cech class gives a global
middle translation.  Its overlap equalities are exactly the two
finite formulas and the cusp formula just displayed.  Consequently
the local biholomorphisms glue, with their inverses, over $B$.

For necessity the base is intrinsic even when its anticanonical
pencil changes.  If $F:X\to X'$ is a biholomorphism, the pullback
of the Fubini--Study form by $f'F$ is exact because
$H^2(S^6,\mathbb R)=0$.  Its restriction to a regular complex torus,
wedged with a flat K\"ahler form, has zero integral.  Positivity
forces $f'F$ to be constant on that torus.  The proper analytic
image of $(f,f'F)$ is therefore a curve of degree one over $B$;
normalization and the inverse argument give a base automorphism.
Its three monodromy orders $3,4,\infty$ are distinct, so it fixes
all three marked points and is the identity.  Multiplicities agree.
On ordinary fibres the lift is affine.  The simultaneous rational
centralizer of $A_1,A_2$ is $\mathbb Q I+\mathbb Q E$, as follows
by solving their entrywise commuting equations.  Comparison of the
constant last two columns of $\Pi_c$ and unimodularity then gives
fibre-linear part $\alpha I$, integer difference $h$, and lattice
map $\alpha S_h^{-1}$.  The normalized root base change is
functorial, so the lifted map commutes with the actual deck
generator.  Summing its defect around that generator gives the
two norm conditions above.  A cusp map has an integral toric
winding: at a generic point of a toric component, its ratio of
characters is $t_c^q$ times a holomorphic unit.  Its class in
$\overline\Lambda$ vanishes.  The middle cocycle relation then
gives $p'=\alpha p$.  This proves all the necessary conditions.

Substitute \ref{push:src:E.1.1}.  In the second congruence the extra term
$N'_1-\alpha N_1$ is divisible by $4l$, and hence drops out
modulo $2l$.  We obtain exactly the statement's two congruences.
For fixed data $b\ne0$ excludes $\alpha=-1$.
The invertibility of $N_0$ modulo $k$, of $N_1$ modulo $2l$, and
$\gcd(k,2l)=1$ give $h\in2kl\mathbb Z$.
For varying data the Chinese remainder theorem gives exactly one
shift class modulo $2kl$, and all residue classes of
$\rho_0\bmod k$, $\rho_1\bmod2l$ occur without changing $b$ or
freeness.  Finally a positive $u$-circuit sends $c$ to $c-1$.
Its translation on $\mathbb Z/(2kl)$ is a generator.
\end{proof}

The preceding theorem answers Question 3 of
\cite[Section 1.1]{ViaclovskyLog2026}.  The original return asserted the
same arithmetic classification in \cref{fam:src:B.3.1}; the contribution here is
the explicit dictionary and the complete enlarged-root and collar
argument.

\begin{corollary}[holomorphic dependence of the comparisons]\label{push:src:E.1.3}
For fixed discrete data in \ref{push:src:E.1.2} and a fixed admissible shift $h$,
the constructed comparisons are holomorphic isomorphisms of families
over a sufficiently small parameter disc.  In particular, for sphere
data with multiplicities $(3,4)$ the identification with the standard
family uses the parameter
\[
 W=e^{-\pi i(c+s_2)}=(-1)^{s_2}w.
\]
The smooth monodromy conjugacy class around the positive parameter
circle is unchanged by this identification of twist data.
\end{corollary}
\Uses{the hypotheses, notation and definitions stated here; \cref{push:src:E.1.2}; \cref{push:src:E.3.2}}

\begin{proof}
In \ref{push:src:E.1.2} the rational vectors $b_j$, integral defects $\ell_j$,
and toric exponents are independent of $c$.  All period matrices,
root comparisons, and cusp unit factors are holomorphic in $c$.
On a small Stein parameter disc $D$, the correction sheaf is
$\mathcal V_B\boxtimes\mathcal O_D$.  Its first cohomology on
$B\times D$ is zero: project to $D$, use the same two-chart
complex and $H^1(B,\mathcal V_B)=0$, and then Stein acyclicity.
Thus the affine correction can be chosen holomorphically in $c$.
For $(k,l)=(1,1)$ the only nonvacuous finite congruence is
$h\equiv s_2-s_2'\pmod2$, because the admissible $r_2'$ is odd.
Taking $s_2'=0$ gives the displayed parameter identification.
Two choices of these family comparisons differ fibrewise by an
automorphism of the standard member.  The full automorphism
calculation audited in \ref{push:src:E.3.2} identifies that group with $\mathbb C^*$;
it is connected.  On overlapping parameter arcs the transitions
are therefore isotopic to the identity.  Composing around a circle
changes the resulting diffeomorphism only by such an isotopy and
by conjugation at the chosen base fibre.  Multiplication of the
parameter by $(-1)^{s_2}$ preserves the positive orientation.
\end{proof}

